\subsection{好方向是否进入了候选池？}
为了区分“评分不准”和“没有送来好候选”，我在已完成线上决定以后，离线枚举固定154方向字典中的全部初次一换一机会（initial one-exchange opportunity），再走最多三个按完整收益选择的局部动作。这个teacher只作诊断，没有回传到线上规则。96个可用单元完整，之前端点失败对应的3个单元仍未运行。

在十个完整额外对象中，如果允许anchor在全字典排序，它保留的初次正收益机会为97.30--99.99\%，中位99.54\%。但是固定12个移入候选的短池只覆盖49.12--95.31\%，中位72.95\%。这些比例先对对象内机会求和，因此我把目前更明显的瓶颈理解为候选取得（candidate acquisition），而不是排序公式本身。这并不表示已经解决全局选模：后续交换进入不同条件空间，局部路径也会不同。完整收益贪心路径在90个单元终态更好，线上路径在5个单元更好，1个在评价底噪内相同；这个差距包含覆盖、排序与路径共同作用。
